\documentclass{../templateWriting/cdwriting}

\newcommand{\authorinfo}{THE FORUM CENTER - NGUYỄN HOÀNG HUY}
\newcommand{\documenttitle}{TÀI LIỆU CHUYÊN ĐỀ WRITING}
\newcommand{\collectiontitle}{Level 1 — W6 (Feb 2-8)}

\begin{document}

\thispagestyle{firstpage}
\makeworkshopheader{\authorinfo}{\documenttitle}{\collectiontitle}
\pagestyle{otherpages}

\workshopprompt{Government should invest more money in science education rather than other subjects to develop the country. Do you agree or disagree?}

\cdsection{Phần I: Chiến lược tư duy "Một phe" \textit{(One-sided Strategy)}}

\cdstep{1. Bước 1: Xác định quan điểm \textit{(Positioning)}}
\begin{itemize}
\item \cdkv{Câu hỏi}{Bạn có đồng ý rằng nên dồn tiền cho Khoa học và coi nhẹ các môn khác không?}
\item \cdkv{Câu trả lời}{Tôi \cdgreen{\textbf{KHÔNG ĐỒNG Ý}} (I disagree).}
\item \cdkv{Tư duy chủ đạo}{Tôi sẽ chứng minh rằng các môn xã hội (Văn, Sử, Địa, Nhạc...) là nền tảng của đất nước, nếu thiếu chúng đất nước sẽ gặp nguy hiểm. Tôi sẽ không khen khoa học trong bài này để tránh bị lạc đề sang "both views".}
\end{itemize}

\cdstep{2. Bước 2: Tìm ý \textit{(Brainstorming)} cho phe DISAGREE}

Chúng ta cần 2 lý do lớn để bảo vệ các môn học "Other subjects" (Các môn không phải khoa học).
\begin{itemize}
\item \cdkv{Lý do 1}{Về Văn hóa \& Bản sắc (Culture \& Identity)}
\begin{itemize}
\item \cdpurple{\underline{\textbf{Môn học:}}} Lịch sử (History), Văn học (Literature).
\item \cdpurple{\underline{\textit{Tại sao quan trọng:}}} Nếu chỉ học Toán/Lý, học sinh sẽ quên mất mình là ai, quên lịch sử đất nước. Đất nước sẽ mất đi bản sắc.
\end{itemize}
\item \cdkv{Lý do 2}{Về Tinh thần \& Xã hội (Mental Health \& Society)}
\begin{itemize}
\item \cdpurple{\underline{\textbf{Môn học:}}} Âm nhạc (Music), Mỹ thuật (Arts).
\item \cdpurple{\underline{\textit{Tại sao quan trọng:}}} Khoa học làm ra máy móc, nhưng Nghệ thuật làm người ta vui vẻ, giảm stress. Một đất nước phát triển cần người dân hạnh phúc, không phải những con robot.
\end{itemize}
\end{itemize}

\cdsection{Phần II: Dàn ý chi tiết \textit{(Outline)}}
\begin{itemize}
\item \cdblue{\textbf{Mở bài (Introduction):}}
\begin{itemize}
\item \cdkv{Paraphrase}{Một số người muốn chính phủ ưu tiên tiền cho khoa học.}
\item \cdkv{Thesis Statement (Quan điểm)}{Tôi hoàn toàn không đồng ý. Tôi tin rằng các môn học khác cũng quan trọng và cần được đầu tư tiền.}
\end{itemize}
\item \cdkv{Thân bài 1 (Body 1)}{Tầm quan trọng của Lịch sử và Văn học.}
\begin{itemize}
\item \cdkv{Câu chủ đề}{Các môn xã hội giúp giữ gìn văn hóa.}
\item \cdkv{Giải thích}{Lịch sử dạy về quá khứ. Văn học dạy về ngôn ngữ.}
\item \cdkv{Kết quả}{Giúp người dân yêu nước hơn.}
\end{itemize}
\item \cdkv{Thân bài 2 (Body 2)}{Tầm quan trọng của Nghệ thuật và Âm nhạc.}
\begin{itemize}
\item \cdkv{Câu chủ đề}{Các môn nghệ thuật giúp cải thiện đời sống tinh thần.}
\item \cdkv{Giải thích}{Cuộc sống hiện đại rất áp lực. Âm nhạc/Vẽ giúp thư giãn.}
\item \cdkv{Kết quả}{Xã hội sẽ hạnh phúc và sáng tạo hơn.}
\end{itemize}
\item \cdblue{\textbf{Kết bài (Conclusion):}}
\begin{itemize}
\item \cdkv{Khẳng định lại}{Tôi không đồng ý việc chỉ ưu tiên khoa học. Chính phủ phải đầu tư cho các môn xã hội và nghệ thuật nữa.}
\end{itemize}
\end{itemize}

\cdsection{Phần III: Bài viết mẫu \textit{(Sample Essay)}}

\cdkv{Lưu ý}{Ngữ pháp sử dụng câu đơn, rõ ràng, tập trung hoàn toàn vào việc bảo vệ "Other subjects".}

\cdstep{1. Mở bài}

Some people think that the government should spend more money on science education rather than on other subjects. I totally disagree with this opinion. I believe that subjects like history and arts are extremely important for the development of a country.

\cdstep{2. Thân bài 1 (Lý do 1: Văn hóa \& Lịch sử)}

First of all, subjects like history and literature are essential for our culture.

(Đầu tiên, các môn như lịch sử và văn học rất thiết yếu cho văn hóa của chúng ta.)

History teaches students about the past and the traditions of their country. If the government does not invest in these subjects, young people may forget their origins. For example, knowing about national heroes makes citizens love their country more. Therefore, spending money on these subjects is necessary to keep the national identity.

(Lịch sử dạy học sinh về quá khứ và truyền thống đất nước. Nếu chính phủ không đầu tư vào các môn này, người trẻ có thể quên nguồn gốc. Ví dụ, biết về các anh hùng dân tộc giúp công dân yêu nước hơn. Do đó, chi tiền cho các môn này là cần thiết để giữ gìn bản sắc quốc gia.)

\cdstep{3. Thân bài 2 (Lý do 2: Nghệ thuật \& Tinh thần)}

Secondly, arts and music play a big role in our mental health.

(Thứ hai, nghệ thuật và âm nhạc đóng vai trò lớn trong sức khỏe tinh thần của chúng ta.)

Science creates technology, but arts help people relax. In modern life, people have a lot of stress. Learning music or painting helps students feel happy and creative. A developed country needs happy people, not just machines. So, the government should support art education as well.

(Khoa học tạo ra công nghệ, nhưng nghệ thuật giúp con người thư giãn. Trong cuộc sống hiện đại, mọi người chịu nhiều áp lực. Học nhạc hoặc vẽ giúp học sinh cảm thấy vui vẻ và sáng tạo. Một đất nước phát triển cần những con người hạnh phúc, không chỉ là máy móc. Vì vậy, chính phủ cũng nên hỗ trợ giáo dục nghệ thuật.)

\cdstep{4. Kết bài}

In conclusion, I do not agree that science is more important than other subjects. History, literature, and arts bring great benefits to society. The government should pay attention to these subjects to develop a better country.

(Tóm lại, tôi không đồng ý rằng khoa học quan trọng hơn các môn khác. Lịch sử, văn học và nghệ thuật mang lại lợi ích to lớn cho xã hội. Chính phủ nên chú ý đến các môn học này để phát triển một đất nước tốt đẹp hơn.)

\cdsection{Phần IV: Từ vựng \& Cấu trúc ghi điểm}

\cdstep{1. Nhóm từ thay thế cho "Spend money" (Chi tiền/Đầu tư)}

Thay vì cứ mãi dùng "spend money on" hay "give money to", hãy dùng:
\begin{itemize}
\item \cdkv{Allocate funds to (v)}{Phân bổ ngân sách cho...}
\begin{itemize}
\item \cdkv{Ví dụ}{The government should allocate funds to history and arts.}
\item (Chính phủ nên phân bổ ngân sách cho lịch sử và nghệ thuật.)
\end{itemize}
\item \cdkv{Invest in (v)}{Đầu tư vào... (Từ này cơ bản nhưng rất đúng trọng tâm).}
\begin{itemize}
\item \cdkv{Ví dụ}{It is wise to invest in art education.}
\end{itemize}
\item \cdkv{Financial support (n)}{Sự hỗ trợ tài chính.}
\begin{itemize}
\item \cdkv{Ví dụ}{Arts subjects need more financial support from the state.}
\item (Các môn nghệ thuật cần nhiều sự hỗ trợ tài chính hơn từ nhà nước.)
\end{itemize}
\end{itemize}

\cdstep{2. Nhóm từ thay thế cho "Important" (Quan trọng)}

Trong bài viết 1 phe, bạn sẽ phải khen môn Văn/Sử/Nhạc rất nhiều. Đừng chỉ dùng "important".
\begin{itemize}
\item \cdkv{Crucial (adj)}{Cốt yếu, cực kỳ quan trọng.}
\begin{itemize}
\item \cdkv{Ví dụ}{History is crucial for a nation’s identity.}
\end{itemize}
\item \cdkv{Vital (adj)}{Sống còn, thiết yếu.}
\begin{itemize}
\item \cdkv{Ví dụ}{Music plays a vital role in mental health.}
\end{itemize}
\item \cdkv{Indispensable (adj)}{Không thể thiếu được.}
\begin{itemize}
\item \cdkv{Ví dụ}{Literature is an indispensable part of education.}
\item (Văn học là một phần không thể thiếu của giáo dục.)
\end{itemize}
\end{itemize}

\cdstep{3. Nhóm từ vựng cho Thân bài 1: Văn hóa \& Lịch sử \textit{(Culture \& History)}}

\cdblue{\textbf{Để nói sâu hơn về việc "giữ gìn đất nước":}}
\begin{itemize}
\item \cdkv{Preserve (v)}{Bảo tồn, gìn giữ.}
\begin{itemize}
\item \cdkv{Ví dụ}{We need to preserve our traditional values.}
\item (Chúng ta cần bảo tồn các giá trị truyền thống.)
\end{itemize}
\item \cdkv{Cultural heritage (n)}{Di sản văn hóa.}
\begin{itemize}
\item \cdkv{Ví dụ}{History helps students understand their cultural heritage.}
\end{itemize}
\item \cdkv{Patriotism (n)}{Lòng yêu nước.}
\begin{itemize}
\item \cdkv{Ví dụ}{Learning history can build a sense of patriotism.}
\item (Học lịch sử có thể xây dựng lòng yêu nước.)
\end{itemize}
\item \cdkv{Roots (n)}{Cội nguồn/Gốc rễ.}
\begin{itemize}
\item \cdkv{Ví dụ}{Young people should not forget their roots.}
\end{itemize}
\end{itemize}

\cdstep{4. Nhóm từ vựng cho Thân bài 2: Tinh thần \& Nghệ thuật \textit{(Mental Health \& Arts)}}

Để nói sâu hơn về việc "giúp con người vui vẻ/sáng tạo":
\begin{itemize}
\item \cdkv{Relieve stress (v)}{Giảm căng thẳng (Hay hơn "reduce stress").}
\begin{itemize}
\item \cdkv{Ví dụ}{Painting is a great way to relieve stress.}
\end{itemize}
\item \cdkv{Well-being (n)}{Sự hạnh phúc, khỏe mạnh (về tinh thần).}
\begin{itemize}
\item \cdkv{Ví dụ}{Arts contribute to the well-being of citizens.}
\item (Nghệ thuật đóng góp vào sự hạnh phúc của công dân.)
\end{itemize}
\item \cdkv{Foster creativity (v)}{Nuôi dưỡng sự sáng tạo.}
\begin{itemize}
\item \cdkv{Ví dụ}{Music lessons help foster creativity in children.}
\end{itemize}
\item \cdkv{Holistic development (n)}{Sự phát triển toàn diện (Cụm từ "ăn điểm" cao).}
\begin{itemize}
\item \cdkv{Nghĩa là}{Phát triển cả trí tuệ (khoa học) lẫn tâm hồn (nghệ thuật).}
\item \cdkv{Ví dụ}{For holistic development, students need to learn both science and arts.}
\end{itemize}
\end{itemize}
\cdstep{Bảng nâng cấp từ vựng (Before \& After)}
\cdnoteinline{Dưới đây là cách bạn áp dụng từ mới vào bài viết cũ để thấy sự khác biệt:}
\begin{cdtablebox}
\renewcommand{\arraystretch}{1.25}
\begin{tabularx}{\linewidth}{>{\raggedright\arraybackslash}X>{\raggedright\arraybackslash}X}
\rowcolor{topicblue!12}\textbf{Câu gốc (Band 4.0)} & \textbf{Câu nâng cấp (Band 5.5+)} \\
\hline
The government should spend money on arts. & The government should allocate funds to arts. \\
\hline
History is very important for culture. & History is crucial for preserving cultural heritage. \\
\hline
If we don't learn history, we forget where we come from. & If we ignore history, we may forget our roots. \\
\hline
Arts help people relax. & Arts help people relieve stress and improve well-being. \\
\hline
Music makes students creative. & Music helps foster creativity in students. \\
\hline
\end{tabularx}
\end{cdtablebox}

\end{document}
